\section{An explicit nonabelian arithmetic realization}
\label{sec:nonabelian-frame}

The local theta machinery has more scope than its cyclotomic application
suggests. Its multivariable formal series allow an arbitrary permutation
of the embedding coordinates. The arithmetic specialization identifies
that permutation with multiplication by $-p^e$ on cyclotomic embeddings.
We construct a concrete replacement involving a totally real $S_3$
field: an exact integral frame, the corresponding local Frobenius
congruence, and a selector which distinguishes complete splitting from
an irreducible cubic at primes in the same arithmetic progression.

These are proved arithmetic inputs. They do not supply a new global
norm inequality or prove nonabelian GRH. Their role is to identify a
specific transfer which is feasible locally and the exact global
statements still needed.

\subsection{The field and its compatible embedding set}

Put
\[
 f(X)=X^3-4X+1,\qquad F_0=\Q(\alpha),\quad f(\alpha)=0,
 \qquad F=\text{the splitting field of }f.
\]
Neither $1$ nor $-1$ is a root, so $f$ is irreducible. It has a real
root in each of $(-3,-2)$, $(0,1)$, and $(1,2)$, and its discriminant is
\[
 -4(-4)^3-27=229.
\]
Consequently $F$ is totally real,
$\Gal(F/\Q)\simeq S_3$, and $F=F_0(\sqrt{229})$.
Only 229 can ramify: the polynomial has discriminant 229, and its
splitting field is unramified away from that discriminant.

Take source-compatible parameters with
$18480\cdot229\mid q_0$, and a cyclotomic tower level $q_1$ divisible
by $q_0$. Let $K=\Q(\sqrt{-d})$, with $d$ and $q_0$ satisfying the
source's admissibility conditions. Define
\[
 E=KF\Q(\zeta_{q_1}),\qquad B_0=F\Q(\zeta_{q_0}).
\]
The source's splitting condition makes every prime dividing $q_1$
split in $K$; see
\src{New/Reduction/Data/FieldsAndPlaces}{chiK_eq_one_of_dvd},
line 144. In particular, every prime inert in $K$ is unramified in
$F\Q(\zeta_{q_1})$.

The quadratic intersection must be retained:
\begin{equation}\label{eq:s3-intersection}
 F\cap\Q(\zeta_{q_1})=\Q(\sqrt{229}).
\end{equation}
Indeed the intersection is an abelian subextension of the $S_3$
extension, hence lies in its unique quadratic subfield; that subfield
lies in $\Q(\zeta_{229})$. Writing $\chi_{229}$ for the quadratic
character of conductor 229 gives
\begin{equation}\label{eq:s3-fiber-product}
 \Gal(F\Q(\zeta_{q_1})/\Q)
 \simeq
 \left\{(g,u)\in S_3\times(\Z/q_1\Z)^\times:
 \operatorname{sgn}(g)=\chi_{229}(u)\right\}.
\end{equation}
Its order is $3\varphi(q_1)$. The field $K$ is disjoint from this
compositum: it ramifies at $d$, where the compositum is unramified.
Thus the induced CM type of $E$ has $3\varphi(q_1)$ embeddings,
and $\Gal(E/K)$ is the fiber product in
\eqref{eq:s3-fiber-product}. Replacing it by a direct product would
give both the wrong dimension and the wrong prime-splitting density.

\subsection{An exact integral frame over the real $S_3$ field}

\begin{theorem}[Explicit arithmetic frame]\label{thm:s3-frame}
Set
\[
 u_1=1,\qquad u_2=\alpha,\qquad
 u_3=24\alpha^2+9\alpha-64,
\]
and let
\[
 (v_1,\ldots,v_6)
 =(u_1,u_2,u_3,\sqrt{229}u_1,\sqrt{229}u_2,\sqrt{229}u_3).
\]
With
\[
 M_F=2517168=48\cdot229^2,
 \qquad
 (m_1,\ldots,m_6)=(419528,157323,229,1832,687,1),
\]
all $v_i$ are algebraic integers and, for any two real embeddings
$\sigma,\tau:F\hookrightarrow\R$,
\begin{equation}\label{eq:s3-frame}
 \sum_{i=1}^{6}m_i\sigma(v_i)\tau(v_i)
 =M_F\delta_{\sigma\tau}.
\end{equation}
\end{theorem}

\begin{proof}
Newton sums for $f$ give
\[
 \Tr_{F_0/\Q}(1,\alpha,\alpha^2,\alpha^3,\alpha^4)
 =(3,0,8,-3,32).
\]
The trace Gram matrix in $(1,\alpha,\alpha^2)$ and the change-of-basis
matrix with columns $u_i$ are
\[
 T=\begin{pmatrix}3&0&8\\0&8&-3\\8&-3&32\end{pmatrix},
 \qquad
 U=\begin{pmatrix}1&0&-64\\0&1&9\\0&0&24\end{pmatrix}.
\]
Exact multiplication yields
\begin{equation}\label{eq:s3-trace-orthogonal}
 U^{\mathsf T}TU=\operatorname{diag}(3,8,5496).
\end{equation}
The nontrivial automorphism of $F/F_0$ negates $\sqrt{229}$.
The cross terms between $F_0$ and $\sqrt{229}F_0$ therefore have
trace zero, and the six $v_i$ form a trace-orthogonal basis with
squared trace norms
\begin{equation}\label{eq:s3-norms}
 (d_1,\ldots,d_6)=(6,16,10992,1374,3664,2517168).
\end{equation}

Let $V=(\sigma(v_i))_{\sigma,i}$ be the square embedding matrix.
Then $V^{\mathsf T}V=\operatorname{diag}(d_i)$, and $V$ is
invertible. Hence
\[
 V\operatorname{diag}(d_i^{-1})V^{\mathsf T}=I.
\]
The listed positive integers satisfy $m_i=M_F/d_i$.
Multiplying by $M_F$ proves \eqref{eq:s3-frame}.
All basis elements are integral because $\alpha$ and $\sqrt{229}$
are integral.
\end{proof}

For the cubic field alone, the same proof gives the smaller frame
\begin{equation}\label{eq:cubic-frame}
 1832\sigma(u_1)\tau(u_1)+687\sigma(u_2)\tau(u_2)
 +\sigma(u_3)\tau(u_3)=5496\delta_{\sigma\tau}.
\end{equation}
The $S_3$ closure is used here because its embeddings form a group,
which permits the coset selector below.

\begin{proposition}[Extension across the cyclotomic intersection]
\label{prop:s3-tensor}
The elements $w_{ij}=v_i\zeta_{q_1}^{j}$, for
$1\le i\le6$ and $0\le j<q_1$, with multiplicity $m_i$, form a
tight frame of bound
\[
 B=M_Fq_1
\]
on the induced CM type of $E$. More precisely, for its embeddings
$\rho,\rho'$,
\begin{equation}\label{eq:s3-tensor-frame}
 \sum_{i,j}m_i\rho(w_{ij})\overline{\rho'(w_{ij})}
 =B\delta_{\rho\rho'}.
\end{equation}
The vectors $w_{ij}/B$, with multiplicities $Bm_i$, give a Parseval
resolution. At every prime inert in $K$ these normalized vectors are
integral above that prime.
\end{proposition}

\begin{proof}
An embedding of $F\Q(\zeta_{q_1})$ is a pair $(\sigma,\eta)$ of
embeddings of its two factors which agree on $\Q(\sqrt{229})$.
The cyclotomic geometric sum is
\[
 \sum_{j=0}^{q_1-1}\eta(\zeta_{q_1}^{j})
 \overline{\eta'(\zeta_{q_1}^{j})}
 =q_1\delta_{\eta\eta'}.
\]
Multiplying by \eqref{eq:s3-frame} gives
$B\delta_{\sigma\sigma'}\delta_{\eta\eta'}$.
Restricting to compatible pairs gives
\eqref{eq:s3-tensor-frame}, since the two restrictions determine
the embedding of the compositum. This argument uses no linear
disjointness assumption.

Dividing each vector by $B$ divides the frame operator by $B^2$;
multiplying each multiplicity by $B$ then makes it the identity.
Equivalently, replace each occurrence of a frame vector by $B$ copies
in total ($B-1$ additional copies), before dividing by $B$.
Finally every prime factor of $B=48\cdot229^2q_1$ is split in $K$,
whereas each $w_{ij}$ is integral. For an inert rational prime $p$,
the denominator $B$ is a unit at every place above $p$. Thus all
normalized coefficients belong to the $p$-localized ring
$\mathcal O_E\otimes_{\Z}\Z_{(p)}$; the same is true of their
conjugates.
\end{proof}

Placing one copy of this resolution in each coordinate gives the
identity on $E^n$ for any finite $n$. This is the algebraic
\texttt{IsResolution} input in
\src{New/Derivatives/FormalGroup/Resolution}{IsResolution},
lines 48--92, and the normalization agrees with the construction in
\src{New/Derivatives/Vacuum/StandardDataProducts}{resolutionVectors},
lines 62--100.

\begin{remark}[What the frame does not establish]
The source's global \texttt{IsTightFrame} additionally places the
vectors in a specified Weil lattice $E\cap\alpha_{\rm Weil}\widehat
{\mathcal O}_E$. We have proved the embedding identity and the local
integrality at inert-$K$ primes. We have not constructed the required
global Weil data, a frame of exactly the principal level $p_0p_1$,
two coprime frame bounds, or a frame member satisfying the source's
\texttt{HasUnitLinearForm} condition at every prime. These distinctions
are used in the global proof.
\end{remark}

\subsection{The actual local nonabelian partner}

\begin{proposition}[Frobenius partner at inert primes]
\label{prop:s3-partner}
Let $p$ be inert in $K$, let $\iota:E\hookrightarrow\C_p$ select
a prime $\mathfrak p$ of $E$, and let $h\in\Gal(E/\Q)$ be
arithmetic Frobenius at $\mathfrak p$. Let $c$ be complex conjugation
and identify the induced CM type with the embeddings restricting
to the chosen embedding of $K$. For odd positive $e$, put
\[
 \operatorname{pr}_{p,e}(\tau)=c\,h^e\tau.
\]
This is again a member of the CM type. For every $a\in\mathcal O_E$,
\begin{equation}\label{eq:s3-frobenius-congruence}
 \iota(\tau(a))^{p^e}
 \equiv\iota\!\left(\overline{\operatorname{pr}_{p,e}(\tau)(a)}\right)
 \pmod {p\mathcal O_{E_{\mathfrak p}}}.
\end{equation}
On the $F$ and cyclotomic components the partner acts respectively
by $\operatorname{Frob}_{F,p}^{e}$ and $-p^e$.
\end{proposition}

\begin{proof}
The prime $p$ is unramified in all factors of $E$. In the unramified
completion, arithmetic Frobenius satisfies
\[
 \iota(h(b))\equiv\iota(b)^p\pmod{p\mathcal O_{E_{\mathfrak p}}}
 \qquad(b\in\mathcal O_E).
\]
Iteration and $b=\tau(a)$ prove
$\iota(h^e\tau(a))\equiv\iota(\tau(a))^{p^e}$.
Conjugation of $c h^e\tau(a)$ is $h^e\tau(a)$, giving
\eqref{eq:s3-frobenius-congruence}. Both $c$ and $h^e$ conjugate
$K$ because $p$ is inert and $e$ is odd; their product fixes its
chosen embedding. Complex conjugation fixes the totally real
field $F$ and inverts $\zeta_{q_1}$, giving the asserted component
actions.
\end{proof}

Unramifiedness here is essential. It turns the Frobenius congruence
into divisibility by $p$, with valuation at least one, which is the
strength used in the source. A congruence only modulo the maximal
ideal of an arbitrarily ramified completion would be weaker.
Changing the $p$-adic embedding conjugates Frobenius. The
identity-class selector below is invariant under that ambiguity.
The separate source calculation at the prime $p=d$, ramified in
$K$, uses reduction modulo $\sqrt{-d}$; it is not included in
Proposition~\ref{prop:s3-partner}.

To see the local connection explicitly, let $u_j\in E^n$ be a
Parseval resolution integral at every place above $p$, with
multiplicities $a_j$,
and put
\[
 L_j(z)=\sum_{\tau,r}\overline{\tau(u_{j,r})}\,z_{\tau,r}.
\]
All coefficients lie in $\mathcal O_E\otimes_{\Z}\Z_{(p)}$, so a
common rational denominator prime to $p$ clears them and their
conjugates. Applying
\eqref{eq:s3-frobenius-congruence} to conjugated coefficients and
using the Parseval identity yields
\begin{equation}\label{eq:s3-coupling}
 \sum_j a_j L_j(z)^{p^e}L_j(z)
 \equiv\sum_{\tau,r}z_{\tau,r}^{p^e}
 z_{\operatorname{pr}_{p,e}(\tau),r}\pmod p.
\end{equation}
For example, the coefficient of $z_{\tau,r}^{p^e}z_{\sigma,s}$
on the left is congruent to
$\sum_j a_j\operatorname{pr}_{p,e}(\tau)(u_{j,r})
\overline{\sigma(u_{j,s})}$, which is the required Kronecker delta.

The generic source lemma
\src{New/Derivatives/FormalGroup/ThetaProducts}{dvd_sum_subst_sub_couplingSum},
lines 301--347, accepts precisely such a coupling congruence for an
arbitrary partner map. Likewise
\src{New/Derivatives/Divisibility/NormalForm}{HasNormalForm},
lines 38--54, imposes no cyclotomic restriction on that map.
The specialized bridge is
\src{New/Derivatives/FormalGroup/ThetaProducts}{norm_pow_sub_conj_partner_le},
lines 138--176, followed by \texttt{exists\_hasNormalForm},
lines 443--480. Proposition~\ref{prop:s3-partner} and
\eqref{eq:s3-coupling} explain how to replace its arithmetic input
at inert primes. The coefficients of the linear forms carry
nonabelian information even though the one-variable theta function
still comes from a CM elliptic curve over $K$.

\subsection{A selector with nonabelian arithmetic content}

\begin{lemma}[Coset correlation]\label{lem:s3-coset}
Let $G$ be a finite group, $H\le G$, $t\in G$, $S\subseteq H$,
and $B=tS$. Then
\begin{equation}\label{eq:s3-coset}
 \#\{a\in H:ga\in B\}=|B|\,\mathbf1_{gH=tH}
 \qquad(g\in G).
\end{equation}
\end{lemma}

\begin{proof}
Left multiplication by $g$ is a bijection from $H$ to $gH$.
If $gH=tH$, its inverse gives one preimage for each element of $B$.
Otherwise those left cosets are disjoint.
\end{proof}

In our arithmetic setting let
\[
 G=\Gal(E/K),\qquad H=\Gal(E/KB_0),\qquad t=(1,-1).
\]
The fiber product shows
\[
 H=\{(1,u):u\equiv1\pmod{q_0}\},\qquad
 [G:H]=3\varphi(q_0).
\]
The subgroup $H$ is central; its cyclotomic tower coordinates permit
the choice of a sparse $S\subseteq H$. For the partner
$g=c h^e$ of an inert prime, the condition $gH=tH$ is exactly
\begin{equation}\label{eq:s3-split-selector}
 \operatorname{Frob}_{B_0,p}^{e}=1.
\end{equation}
For $e=1$ this is complete splitting in $B_0$, with density
$1/(3\varphi(q_0))$. The same formula retains the necessary
prime-power information for higher $e$.

This gives the precise group-theoretic replacement for
\src{New/Reduction/PrimeSums}{card_filter_mem_bPlaces},
lines 101--108. It does not yet replace all arithmetic estimates
which use the old sets of places.

\begin{remark}[A rigorous failure of the naive pullback]
If $\pi:G\twoheadrightarrow Q$ is a quotient and
$A=\pi^{-1}(A_0)$, $B=\pi^{-1}(B_0')$, then
\[
 \#\{a\in A:ga\in B\}
 =|\ker\pi|\,\#\{a_0\in A_0:\pi(g)a_0\in B_0'\}.
\]
Pulling both old selectors back from the cyclotomic quotient
therefore cannot detect Frobenius in the kernel. The new selector
$H,tS$ avoids this precise obstruction.
\end{remark}

An exact numerical example confirms that the distinction is
arithmetical. Take
\[
 q_0=4231920=18480\cdot229,\qquad d=167160839.
\]
The integer $d$ is prime, $d\equiv7\pmod8$,
$q_0\mid2(d+1)$, and $q_0\nmid d+1$. Both primes
\[
 p_1=38087281,\qquad p_2=59246881
\]
are $1\pmod{q_0}$ and inert in $K$. Yet
\[
 f(X)\equiv(X+16188066)(X-13084869)(X-3103197)\pmod {p_1},
\]
whereas $f$ is irreducible modulo $p_2$. Thus the respective
Frobenius classes in $F$ are the identity and a 3-cycle. The
checker verifies primality by trial division, the first
factorization by multiplication, and the second assertion by
$\gcd(f,X^{p_2}-X)=1$; a cubic without a root is irreducible.
It also checks the coset identity exhaustively in a finite
fiber-product model. These computations accompany the mathematical
proofs, rather than replacing them.

\subsection{Three missing global packages}

\paragraph{1. A noncyclotomic odd-cocycle bound.}
The exact obstruction is
\src{Lib/Cyclotomic/OddSha}{New.exists_length_shaOdd_le},
lines 645--651. Its hypothesis places the entire CM field inside
$\Q(\zeta_N)$, and it bounds the source's specified odd cocycle
modules by the valuations of a nonzero integer
\[
 X\le2^{4(N+t+3)^{27}}h_E.
\]
For a field containing our $S_3$ extension, this hypothesis is
impossible for every $N$: a subfield of an abelian Galois extension
is itself Galois and abelian over $\Q$. The bound is used directly in
\src{New/Congruences/FirstOrder/TraceRelations}{oddClassBound},
lines 74--94, and in the summed denominator-error estimate of
\src{New/Congruences/FitTwistsExist}{isNegligible_padicValNat_oddClassBound},
lines 153--171.

A replacement must address these exact cocycles and give controlled
growth as both the cyclotomic tower and the weights vary. Finiteness
at each fixed field and weight, with an unspecified constant, does
not supply the required error budget. The generic anticyclotomic
twisting theorem takes finite obstruction budgets as input; it
does not automatically establish this separate bound.

\paragraph{2. Globally compatible integral data and level primes.}
The source constructs a purely imaginary generator of the different,
its compatible Weil lattice, and frames of carefully prescribed
bounds. See
\texttt{deltaE} and
\src{New/Derivatives/Field/PrimeDecomposition}{differentIdeal_eq_span_deltaE},
lines 353--399,
\src{New/Derivatives/Vacuum/ShortVectors}{exists_isTightFrame_coprime},
lines 508--535, and
\src{New/Derivatives/Vacuum/ShortVectors}{Data.exists_isTightFrame},
lines 787--813.
Our frame bound $M_Fq_1$ has a factor 229. Multiplying such a bound
cannot produce an arbitrary level $p_0p_1$ prime to 229; the source's
coprime-bound construction therefore needs new input. Nor may a
principal different or a monogenic integral presentation be assumed
for an enlarged CM field without proof. Fractional-ideal and adelic
formulations are plausible ways to organize that work.

The level selection also changes. A congruence $p\equiv1\pmod N$
does not force splitting in $F$, even when $N$ is divisible by its
discriminant. The two primes above illustrate the failure. New
splitting and principal-ideal conditions must be realized inside
the source's quantitative level windows; bare infinitude of
appropriate primes does not suffice. The original windows and
construction are in
\src{New/Reduction/Data/AdmissibleData}{levelPrime}, lines 140--161,
and \src{New/Reduction/Data/Existence}{exists_prime_pair_norm},
lines 294--343 and 870--888.

\paragraph{3. Local crystalline defects and quantitative spanning.}
The declarations
\src{New/Congruences/LocalLoss/Basic}{fixes_emb_of_fixes_rootsOfUnity},
lines 151--165, and
\src{New/Congruences/LocalLoss/CrystallineDefect}{mem_range_of_modEq_one},
lines 294--302, use the cyclotomic nature of the field. Their
analogues require genuine information about the new factor's
inertia and Frobenius. The Kummer-letter and sparse-place argument
in \texttt{eventually\_exists\_lettersReadThick}, lines 217--289 of
the latter file, must be checked again for the new embedding set.
Some spacing conditions may survive restriction to a sparse
subset of the cyclotomic tower, but the local defect estimates
have not been proved here.

The abstract second-moment spreading theorem is available for
general finite unitary actions, and substantial arithmetic moment
machinery is stated for general CM fields. Its application to the
new filling data, integral lattices, and selector still requires
quantitative bounds. In particular the specialized theorem
\src{New/Spanning/SecondMoment/SpanningTheorem}{isSpread_goodPoincare},
line 339, is not inherited merely by tripling the embedding count.

\subsection{What this suggests about further conjectures}

The prime coefficient suggested by the local realization is
\[
 \log p\,(1-\chi_K(p)^e)
 \mathbf1_{\operatorname{Frob}_{B_0,p}^{e}=1}.
\]
It concerns nonabelian prime splitting, so it contains information
absent from the original Dirichlet congruence mask. The source's
main-term cancellation compares remainders at two cutoffs, $x$ and
$x_1$; consequently the class-specific constant cancels. This
explains why keeping the elliptic CM period from $K$ is not, by
itself, an obstruction to the proposed signal. The new different,
measures, and vacuum normalization must nevertheless be compared
before a global estimate can be claimed.

After such an estimate, a further analytic bridge would still be
needed. The abstract Landau barrier is general, but the source's
auxiliary zero-free quadratic-twist selection and residue argument
are specialized to Dirichlet L-functions. Complete-splitting
coefficients are a useful first test because they correspond to a
regular representation and Dedekind zeta functions. General Artin
holomorphy involves additional pole and signed-residue issues
\cite{Cogdell}. The standard $S_3$ Artin representation is already
monomial; this example is not a new case of Artin holomorphy.

The achieved advance is narrower and concrete: the arithmetic
frame, the compatible embedding calculation, and the inert-prime
partner congruence show that the local CM-theta mechanism can carry
nonabelian Frobenius data. The exact points at which the global
argument remains cyclotomic are now explicit. No full replacement
for the main norm inequality is asserted.
